% !TeX root = ../main.tex

\begin{innovations}

本文围绕包围几何上的高质量映射构造开展研究，创新性主要体现在以下四个方面。

第一，提出基于广义环绕数的面向有理参数曲线包围几何的解析查询框架，将复杂曲边边界上的内外判定统一转化为可解析求解的有理积分问题，并进一步处理边界点情形与导数计算，实现了复杂曲边场景下高效、鲁棒的内外判定。

第二，提出壳空间的高阶壳离散化及其优化构造框架，在满足包围、角度和厚度约束的前提下，定义壳空间中的双射投影，获得更光滑、空间分布更均匀的属性传递通道。

第三，建立笼空间上 Cauchy 坐标及其导数的统一解析公式，并将相关留数计算框架扩展至高阶 Green 坐标，使曲边笼变形在保持保角特性的同时具备更光滑的的变形能力和更自然的交互控制方式。

第四，构建面向笼空间的双调和坐标高阶 BEM 框架，在严格边界插值与内部低扭曲之间给出连续可调的统一表达，为复杂曲边笼编辑提供更稳定、更自然的形变传播机制。

上述结果共同构成了从几何谓词、几何构造到解析计算与数值求解的完整技术链路。
\end{innovations}
